% GENERATED FILE. Edit canonical lessons, then run npm run generate:books.
% canonical-source-hash: fnv1a64:78125579895bc004
% canonical-lessons: HI-C502-pota, HI-C502-nati, HI-C502-bhanja, HI-C502-jhulana, HI-C502-iklauta

\chapter{Grandchildren and Nephews}
\label{ch:hi-grandchildren-and-nephews}
\hlchaptermodality{502}

\begin{chapteropening}
\textbf{By the end of this chapter:} \emph{I can name more members of the family.}

Complete the last lesson of chapter 502: I can name more members of the family.
\end{chapteropening}

\section[\texorpdfstring{\hi{पोता} (potā)}{पोता (potā)}]{\hi{पोता} (potā) — a grandson (a son's son)}

\label{lesson:HI-C502-pota}



\begin{quote}
Before the new one: say the Hindi for to keep up, then the Hindi for a reconciliation.
\end{quote}

\subsection*{You'll want to know: \hi{पोता}}
\textbf{\hi{पोता}} — \emph{potā} — \textquotedblleft{}a grandson (a son's son)\textquotedblright{}.

\begin{quote}
\hi{दादा} \hi{अपने} \hi{पोते} \hi{के} \hi{साथ} \hi{खेल} \hi{रहे} \hi{हैं।}
\end{quote}

\begin{quote}
\emph{dādā apne pote ke sāth khel rahe haiṁ.}
\end{quote}

\begin{quote}
Grandfather is playing with his grandson.
\end{quote}

Masculine.

\begin{cousinweb}
\textbf{\hi{पोता}} (\emph{potā}) is a tadbhava: CDIAL lists Hindi \emph{potā}, \textquotedblleft{}son's son\textquotedblright{}, under Sanskrit \emph{paútra}, \textquotedblleft{}grandson\textquotedblright{} (CDIAL 8416). The granddaughter is \emph{potī}.
\end{cousinweb}

\begin{grammarlens}[title={What you've built}]
One more word for describing what you lived and felt.
\end{grammarlens}

\subsection*{Your turn}
\begin{itemize}
\raggedright
  \item \emph{Say it:} \emph{potā}
  \item \emph{Say it:} \emph{potā}, once more
  \item \emph{Say it:} \emph{dādā apne pote ke sāth khel rahe haiṁ.} — Grandfather is playing with his grandson.
  \item \emph{Recall:} say the Hindi for to keep up, then the Hindi for a reconciliation, then say \emph{potā} again
\end{itemize}

\begin{tcolorbox}[breakable,colback=teal!4,colframe=teal!35!black,title={Before you move on}]
What does \hi{पोता} mean? (\textquotedblleft{}A grandson (a son's son)\textquotedblright{}.) Say it once more.
\end{tcolorbox}
\section[\texorpdfstring{\hi{नाती} (nātī)}{नाती (nātī)}]{\hi{नाती} (nātī) — a grandson (a daughter's son)}

\label{lesson:HI-C502-nati}



\begin{quote}
Before the new one: say the Hindi for a reconciliation, then the Hindi for a grandson (a son's son).
\end{quote}

\subsection*{You'll want to know: \hi{नाती}}
\textbf{\hi{नाती}} — \emph{nātī} — \textquotedblleft{}a grandson (a daughter's son)\textquotedblright{}.

\begin{quote}
\hi{नानी} \hi{अपने} \hi{नाती} \hi{को} \hi{कहानी} \hi{सुना} \hi{रही} \hi{हैं।}
\end{quote}

\begin{quote}
\emph{nānī apne nātī ko kahānī sunā rahī haiṁ.}
\end{quote}

\begin{quote}
Grandmother is telling her grandson a story.
\end{quote}

Masculine.

\begin{cousinweb}
\textbf{\hi{नाती}} (\emph{nātī}) is a tadbhava: CDIAL lists Hindi \emph{nātī}, \textquotedblleft{}daughter's son\textquotedblright{}, under Sanskrit \emph{náptṛ}, \textquotedblleft{}grandson\textquotedblright{} (CDIAL 6955b). Hindi keeps the father's side (\emph{potā}) and the mother's side (\emph{nātī}) apart.
\end{cousinweb}

\begin{grammarlens}[title={What you've built}]
One more word for describing what you lived and felt.
\end{grammarlens}

\subsection*{Your turn}
\begin{itemize}
\raggedright
  \item \emph{Say it:} \emph{nātī}
  \item \emph{Say it:} \emph{nātī}, once more
  \item \emph{Say it:} \emph{nānī apne nātī ko kahānī sunā rahī haiṁ.} — Grandmother is telling her grandson a story.
  \item \emph{Recall:} say the Hindi for a reconciliation, then the Hindi for a grandson (a son's son), then say \emph{nātī} again
\end{itemize}

\begin{tcolorbox}[breakable,colback=teal!4,colframe=teal!35!black,title={Before you move on}]
What does \hi{नाती} mean? (\textquotedblleft{}A grandson (a daughter's son)\textquotedblright{}.) Say it once more.
\end{tcolorbox}
\section[\texorpdfstring{\hi{भांजा} (bhāñjā)}{भांजा (bhāñjā)}]{\hi{भांजा} (bhāñjā) — a nephew (a sister's son)}

\label{lesson:HI-C502-bhanja}



\begin{quote}
Before the new one: say the Hindi for a grandson (a son's son), then the Hindi for a grandson (a daughter's son).
\end{quote}

\subsection*{You'll want to know: \hi{भांजा}}
\textbf{\hi{भांजा}} — \emph{bhāñjā} — \textquotedblleft{}a nephew (a sister's son)\textquotedblright{}.

\begin{quote}
\hi{मेरा} \hi{भांजा} \hi{अगले} \hi{साल} \hi{स्कूल} \hi{जाएगा।}
\end{quote}

\begin{quote}
\emph{merā bhāñjā agle sāl skūl jāegā.}
\end{quote}

\begin{quote}
My nephew will start school next year.
\end{quote}

Masculine.

\begin{cousinweb}
\textbf{\hi{भांजा}} (\emph{bhāñjā}) is a tadbhava: CDIAL lists Hindi \emph{bhā̃jā}, \textquotedblleft{}sister's son\textquotedblright{}, under Sanskrit \emph{bhāginēya}, \textquotedblleft{}sister's son\textquotedblright{} (CDIAL 9433), from \emph{bhaginī}, \textquotedblleft{}sister\textquotedblright{}. A brother's son is \emph{bhatījā}.
\end{cousinweb}

\begin{grammarlens}[title={What you've built}]
One more word for describing what you lived and felt.
\end{grammarlens}

\subsection*{Your turn}
\begin{itemize}
\raggedright
  \item \emph{Say it:} \emph{bhāñjā}
  \item \emph{Say it:} \emph{bhāñjā}, once more
  \item \emph{Say it:} \emph{merā bhāñjā agle sāl skūl jāegā.} — My nephew will start school next year.
  \item \emph{Recall:} say the Hindi for a grandson (a son's son), then the Hindi for a grandson (a daughter's son), then say \emph{bhāñjā} again
\end{itemize}

\begin{tcolorbox}[breakable,colback=teal!4,colframe=teal!35!black,title={Before you move on}]
What does \hi{भांजा} mean? (\textquotedblleft{}A nephew (a sister's son)\textquotedblright{}.) Say it once more.
\end{tcolorbox}
\section[\texorpdfstring{\hi{झुलाना} (jhulānā)}{झुलाना (jhulānā)}]{\hi{झुलाना} (jhulānā) — to rock; to swing (someone)}

\label{lesson:HI-C502-jhulana}



\begin{quote}
Before the new one: say the Hindi for a grandson (a daughter's son), then the Hindi for a nephew (a sister's son).
\end{quote}

\subsection*{You'll want to know: \hi{झुलाना}}
\textbf{\hi{झुलाना}} — \emph{jhulānā} — \textquotedblleft{}to rock; to swing (someone)\textquotedblright{}.

\begin{quote}
\hi{दादी} \hi{पोते} \hi{को} \hi{झूले} \hi{में} \hi{झुला} \hi{रही} \hi{हैं।}
\end{quote}

\begin{quote}
\emph{dādī pote ko jhūle meṁ jhulā rahī haiṁ.}
\end{quote}

\begin{quote}
Grandmother is rocking her grandson in the swing.
\end{quote}

Past, with \hi{ने} (he or she …): \hi{उसने} \hi{बच्चे} \hi{को} \hi{झुलाया}, \emph{usne bacce ko jhulāyā}.

\begin{cousinweb}
\textbf{\hi{झुलाना}} (\emph{jhulānā}) is the causative partner of \emph{jhūlnā}, \textquotedblleft{}to swing\textquotedblright{}, which CDIAL lists under \emph{*jhulyati}, a form it reconstructs (CDIAL 5406); \emph{jhūlā}, \textquotedblleft{}a swing\textquotedblright{}, belongs to the same family.
\end{cousinweb}

\begin{grammarlens}[title={What you've built}]
One more word for describing what you lived and felt.
\end{grammarlens}

\subsection*{Your turn}
\begin{itemize}
\raggedright
  \item \emph{Say it:} \emph{jhulānā}
  \item \emph{Say it:} \emph{jhulānā}, once more
  \item \emph{Say it:} \emph{dādī pote ko jhūle meṁ jhulā rahī haiṁ.} — Grandmother is rocking her grandson in the swing.
  \item \emph{Recall:} say the Hindi for a grandson (a daughter's son), then the Hindi for a nephew (a sister's son), then say \emph{jhulānā} again
\end{itemize}

\begin{tcolorbox}[breakable,colback=teal!4,colframe=teal!35!black,title={Before you move on}]
What does \hi{झुलाना} mean? (\textquotedblleft{}To rock; to swing (someone)\textquotedblright{}.) Say it once more.
\end{tcolorbox}
\section[\texorpdfstring{\hi{इकलौता} (iklautā)}{इकलौता (iklautā)}]{\hi{इकलौता} (iklautā) — only (child)}

\label{lesson:HI-C502-iklauta}



\begin{quote}
Before the new one: say the Hindi for a nephew (a sister's son), then the Hindi for to rock.
\end{quote}

\subsection*{You'll want to know: \hi{इकलौता}}
\textbf{\hi{इकलौता}} — \emph{iklautā} — \textquotedblleft{}only (child)\textquotedblright{}.

\begin{quote}
\hi{अरुण} \hi{अपने} \hi{माँ} \hi{बाप} \hi{का} \hi{इकलौता} \hi{बेटा} \hi{है।}
\end{quote}

\begin{quote}
\emph{Aruṇ apne māṁ bāp kā iklautā beṭā hai.}
\end{quote}

\begin{quote}
Arun is his parents' only son.
\end{quote}

\begin{cousinweb}
\textbf{\hi{इकलौता}} (\emph{iklautā}) is a tadbhava: CDIAL lists Hindi \emph{iklautā}, \textquotedblleft{}only begotten, only son\textquotedblright{}, under \emph{*ēkkalaputra}, a form it reconstructs from \emph{ēkkala}, \textquotedblleft{}alone\textquotedblright{}, and \emph{putra}, \textquotedblleft{}son\textquotedblright{} (CDIAL 2507). \emph{akelā}, \textquotedblleft{}alone\textquotedblright{}, comes from the same \emph{ēkkala}.
\end{cousinweb}

\begin{grammarlens}[title={What you've built}]
One more word for describing what you lived and felt.
\end{grammarlens}

\subsection*{Your turn}
\begin{itemize}
\raggedright
  \item \emph{Say it:} \emph{iklautā}
  \item \emph{Say it:} \emph{iklautā}, once more
  \item \emph{Say it:} \emph{Aruṇ apne māṁ bāp kā iklautā beṭā hai.} — Arun is his parents' only son.
  \item \emph{Recall:} say the Hindi for a nephew (a sister's son), then the Hindi for to rock, then say \emph{iklautā} again
\end{itemize}

\begin{tcolorbox}[breakable,colback=teal!4,colframe=teal!35!black,title={Before you move on}]
What does \hi{इकलौता} mean? (\textquotedblleft{}Only (child)\textquotedblright{}.) Say it once more.
\end{tcolorbox}
